% answers_short.tex —— 2009 年·普通高等学校招生考试·北京卷·理科 · 简短答案速查
% 由 tools/make_short_answers.py 自 answers.tex 的 \gksolhead/\gkanswer 提取，
% 并与 2009/2009_数学真题_北京理_答案要点.md 逐题比对后生成，勿手改（改源 MD 后重跑）。

\begin{gknotice}
\textbf{使用说明}：本节只给最终答案便于快速核对，完整推导见解析版。答案为编写者独立推导，\textbf{非官方标准答案}。
\end{gknotice}

\gkpart{选择题}
\begin{center}\begin{tabular}{ccccccccc}
\toprule
题号 & 1 & 2 & 3 & 4 & 5 & 6 & 7 & 8 \\
\midrule
答案 & B & D & C & D & A & C & B & A \\
\bottomrule
\end{tabular}\end{center}

\gkpart{填空题}
\begin{itemize}[leftmargin=2.6em]
  \item[\textbf{9.}] $\frac12$。
  \item[\textbf{10.}] $-6$。
  \item[\textbf{11.}] $-1$。
  \item[\textbf{12.}] $2$；$\frac{2\pi}{3}$。
  \item[\textbf{13.}] $[-3,1]$。
  \item[\textbf{14.}] $1$；$0$。
\end{itemize}

\gkpart{解答题}
\begin{itemize}[leftmargin=2.6em]
  \item[\textbf{15.}] （1）$\sin C=\frac{3+4\sqrt3}{10}$；（2）面积 $\frac{36+9\sqrt3}{50}$。
  \item[\textbf{16.}] （1）证明见解析；（2）$\arcsin\frac{\sqrt2}{4}$；（3）存在，$BD:DP=3:4$。
  \item[\textbf{17.}] （1）$\frac4{27}$；（2）分布列见下，$E\xi=\frac83$。
  \item[\textbf{18.}] （1）$y=x$；（2）见下；（3）$k\in[-1,0)\cup(0,1]$。
  \item[\textbf{19.}] （1）$x^2-\frac{y^2}{2}=1$；（2）$\angle AOB=90^\circ$。
  \item[\textbf{20.}] （1）$\{1,3,4\}$ 不具有，$\{1,2,3,6\}$ 具有；（2）（3）证明见解析。
\end{itemize}
